\documentclass[11pt,a4paper,fontset=fandol]{ctexart}

\usepackage[a4paper,top=2.25cm,bottom=2.45cm,left=2.25cm,right=2.25cm,headheight=18pt,headsep=0.55cm,footskip=24pt]{geometry}
\usepackage{amsmath,amssymb,amsthm}
\usepackage{fancyhdr}
\usepackage{lastpage}
\usepackage{enumitem}
\usepackage{titlesec}
\usepackage{xcolor}
\usepackage{hyperref}
\usepackage{bookmark}
\usepackage[most]{tcolorbox}
\usepackage{setspace}
\usepackage{array}

\hypersetup{
  colorlinks=true,
  linkcolor=blue!50!black,
  urlcolor=blue!50!black
}

\setstretch{1.13}
\setlength{\parindent}{2em}
\setlength{\parskip}{0.25em}
\allowdisplaybreaks
\setlist[enumerate]{itemsep=0.45em, topsep=0.35em, leftmargin=2.4em}
\setlist[itemize]{itemsep=0.25em, topsep=0.25em, leftmargin=1.9em}

\definecolor{mainblue}{RGB}{36,83,140}
\definecolor{lightblue}{RGB}{241,246,252}
\definecolor{lightgray}{RGB}{246,246,246}
\definecolor{rulegray}{RGB}{110,118,128}

\titleformat{\section}
  {\Large\bfseries\color{mainblue}}
  {\thesection.}
  {0.45em}
  {}
  [\vspace{0.2em}\color{rulegray}\titlerule]
\titlespacing*{\section}{0pt}{1.1em}{0.7em}

\pagestyle{fancy}
\fancyhf{}
\fancyhead[L]{\small 函数图像对称周测 B 卷}
\fancyhead[C]{\small 代数专题：综合判断与变式}
\fancyhead[R]{\small 刘文博}
\fancyfoot[C]{\small 第 \thepage 页 / 共 \pageref{LastPage} 页}
\renewcommand{\headrulewidth}{0.55pt}
\renewcommand{\footrulewidth}{0.35pt}

\newtcolorbox{infobox}[1]{
  breakable,
  colback=lightblue,
  colframe=mainblue,
  boxrule=0.7pt,
  arc=1mm,
  title=\textbf{#1},
  fonttitle=\bfseries
}

\newenvironment{solution}{\par\noindent\textbf{参考解答：}}{\hfill$\square$\par}
\newcommand{\answerspace}{\vspace{2.3cm}}
\newcommand{\biganswerspace}{\vspace{3.2cm}}

\begin{document}

\thispagestyle{empty}
\begin{center}
{\fontsize{22}{28}\selectfont\bfseries 函数图像对称周测 B 卷\par}
\vspace{0.4cm}
{\large 适合初中阶段函数专题提高训练\par}
\end{center}

\begin{infobox}{考试说明}
\begin{itemize}
  \item 测试时间：70--75 分钟
  \item 满分：100 分
  \item B 卷更强调方法迁移，请尽量说明“为什么这样对称”，不要只写最后结果。
\end{itemize}
\end{infobox}

\section{试题}

\begin{enumerate}
  \item （12 分）请用“图像上的点跟着变化”的方法说明：为什么函数
  \[
  y=f(x)
  \]
  关于 y 轴对称后，新方程应写成
  \[
  y=f(-x).
  \]
  \biganswerspace

  \item （12 分）把函数
  \[
  y=x^2-4x+3
  \]
  关于 y 轴对称，求化简后的新方程。
  \answerspace

  \item （12 分）已知函数
  \[
  y=|x|
  \]
  经过某次对称后，顶点仍然在原点，且图像开口向下。求新方程，并说明是哪种对称。
  \answerspace

  \item （12 分）已知函数
  \[
  y=f(x)
  \]
  关于原点对称后变成
  \[
  y=-f(-x).
  \]
  如果原图像上有点 $(3,7)$，那么对称后对应点的坐标是什么？
  \answerspace

  \item （12 分）已知函数
  \[
  y=(x+2)^2-5,
  \]
  说出它关于 y 轴对称后的方程，并说明顶点怎样变化。
  \answerspace

  \item （12 分）已知函数
  \[
  y=x^2+6x+8,
  \]
  说出它是由
  \[
  y=x^2-6x+8
  \]
  经过怎样的对称得到的。
  \answerspace

  \item （14 分）将函数
  \[
  y=2x+1
  \]
  先关于 y 轴对称，再作一次竖直平移后，恰好经过原点。求第二次变化的方向、距离以及最终方程。
  \biganswerspace

  \item （14 分）已知原图像上的点 $(-2,5)$ 在函数
  \[
  y=f(x)
  \]
  上。若新图像由原图像先关于 x 轴对称，再关于 y 轴对称得到，求新图像的方程，并说明点的对应变化。
  \biganswerspace
\end{enumerate}

\newpage
\section{详细解答}

\begin{enumerate}
  \item \begin{solution}
  设原图像上有一点 $(x_0,y_0)$，因为它在函数
  \[
  y=f(x)
  \]
  上，所以有
  \[
  y_0=f(x_0).
  \]

  如果图像关于 y 轴对称，那么这个点会变成
  \[
  (-x_0,y_0).
  \]

  现在把新图像上的横坐标记作 $x$，那么它对应原图像上的横坐标应是
  \[
  -x.
  \]
  因此新图像上的函数值应为
  \[
  y=f(-x).
  \]

  所以，函数 $y=f(x)$ 关于 y 轴对称后，新方程应写成
  \[
  \boxed{y=f(-x)}.
  \]
  \end{solution}

  \item \begin{solution}
  关于 y 轴对称，应把整个式子里的 $x$ 换成 $-x$：
  \[
  y=(-x)^2-4(-x)+3=x^2+4x+3.
  \]
  所以新方程为
  \[
  \boxed{y=x^2+4x+3}.
  \]
  \end{solution}

  \item \begin{solution}
  函数 $y=|x|$ 的顶点是 $(0,0)$。

  如果对称后顶点仍在原点，且图像开口向下，那么它应变成倒着开的 V 形，也就是
  \[
  \boxed{y=-|x|}.
  \]

  这可以由\textbf{关于 x 轴对称}得到；由于 $|x|$ 本来关于 y 轴对称，所以先关于 y 轴再关于 x 轴，也就是关于原点对称，结果同样是 $y=-|x|$。
  \end{solution}

  \item \begin{solution}
  关于原点对称时，点 $(x,y)$ 会变成 $(-x,-y)$。

  因此点 $(3,7)$ 对应到
  \[
  \boxed{(-3,-7)}.
  \]
  \end{solution}

  \item \begin{solution}
  关于 y 轴对称，应把 $x$ 换成 $-x$：
  \[
  y=((-x)+2)^2-5=(x-2)^2-5.
  \]
  所以新方程为
  \[
  \boxed{y=(x-2)^2-5}.
  \]

  原函数顶点是 $(-2,-5)$，关于 y 轴对称后变成
  \[
  \boxed{(2,-5)}.
  \]
  \end{solution}

  \item \begin{solution}
  把函数
  \[
  y=x^2-6x+8
  \]
  关于 y 轴对称，应把式子里的 $x$ 换成 $-x$：
  \[
  y=(-x)^2-6(-x)+8=x^2+6x+8.
  \]
  因此题目中的函数是由原函数\textbf{关于 y 轴对称}得到的。
  \end{solution}

  \item \begin{solution}
  第一步先关于 y 轴对称：
  \[
  y=2(-x)+1=-2x+1.
  \]

  现在它还没有经过原点，因为当 $x=0$ 时，
  \[
  y=1.
  \]

  要让它经过原点，就需要再向下平移 $1$ 个单位。

  所以第二次变化是\textbf{向下平移 $1$ 个单位}，最终方程为
  \[
  y=-2x+1-1=-2x.
  \]
  即
  \[
  \boxed{y=-2x}.
  \]
  \end{solution}

  \item \begin{solution}
  先关于 x 轴对称，点 $(-2,5)$ 变成
  \[
  (-2,-5).
  \]
  再关于 y 轴对称，变成
  \[
  (2,-5).
  \]

  因此这个复合变化等价于关于原点对称。

  所以新图像方程为
  \[
  \boxed{y=-f(-x)}.
  \]

  点的对应变化为
  \[
  (-2,5)\longrightarrow (-2,-5)\longrightarrow (2,-5).
  \]
  \end{solution}
\end{enumerate}

\end{document}
